\documentclass[11pt]{article}
\usepackage{ap-psych}
\title{Unit 3: Development and Learning}
\author{Paul Wang}
\date{November 2026}
\docpart{Unit 3 \textperiodcentered\ Study Guide \textperiodcentered\ 15--25\% of MCQ}

\begin{document}
\maketitle

\begin{abstract}
Unit 3 pairs two questions: how people change across the lifespan, and how experience changes behavior. The development half covers themes and methods (3.1), physical development (3.2), gender and sexual orientation (3.3), cognitive development (3.4), language (3.5), and social-emotional development (3.6). The learning half covers classical conditioning (3.7), operant conditioning (3.8), and social, cognitive, and neurological factors in learning (3.9). Expect many items that describe a child's behavior and ask for the stage, or describe a training routine and ask for the conditioning term. Each section ends with check questions; the key is on the last page.
\end{abstract}

\scoringline{Check questions: \blank\ / \qtotal \qquad Target: 20+ \qquad Reinforcement/punishment grid correct on first try? \blank[3em]}

\section{Themes and methods (3.1)}

Three debates frame the field: \textbf{nature vs.\ nurture}; \textbf{continuity vs.\ discontinuity} (gradual change vs.\ distinct stages); \textbf{stability vs.\ change} (do traits persist across life?). Developmental research relies on \textbf{longitudinal} designs (same people over time) and \textbf{cross-sectional} designs (different ages at one time, vulnerable to cohort effects).

\begin{q}{3}{3.1 Themes}{Distinct, qualitatively different stages = discontinuity.}
A theorist argues that children move through qualitatively distinct stages rather than improving gradually. This position emphasizes
\chh{continuity}{stability}{discontinuity}{nurture}
\end{q}

\section{Physical development (3.2)}

\textbf{Prenatal}: \textbf{teratogens} are harmful agents (alcohol, drugs, viruses) that cross the placenta; \textbf{fetal alcohol syndrome} causes facial abnormalities and cognitive deficits. \textbf{Newborn reflexes}: \textbf{rooting} (turning toward a cheek touch), sucking, grasping. \textbf{Maturation}: biological growth that sets the timing of milestones such as walking. \textbf{Critical/sensitive periods}: windows when certain experiences are needed for normal development.

\textbf{Adolescence} begins with \textbf{puberty}: \textbf{primary sex characteristics} (reproductive organs) and \textbf{secondary sex characteristics} (breasts, voice change, body hair); \textbf{menarche} (first menstruation) and \textbf{spermarche} (first ejaculation). The \textbf{prefrontal cortex} matures into the mid-20s, after the limbic system, which helps explain adolescent risk-taking.

\textbf{Adulthood}: physical abilities peak in the 20s; \textbf{menopause} ends menstruation around age 50. Aging slows reaction time and sensory acuity, but exercise and mental activity protect brain function.

\begin{q}{1}{3.2 Prenatal}{Alcohol crossing the placenta is a teratogen; FAS is the outcome.}
A child born with a small head, distinctive facial features, and learning difficulties after heavy prenatal alcohol exposure most likely has
\ch{fetal alcohol syndrome, caused by a teratogen}{a genetic disorder unrelated to environment}{delayed maturation of reflexes}{a sensitive-period language delay}
\end{q}

\begin{q}{4}{3.2 Adolescence}{Late-maturing prefrontal cortex + earlier-maturing limbic reward system explains risk-taking.}
Which best explains why adolescents take more risks than adults?
\ch{Their cerebellum is fully mature.}{They lack secondary sex characteristics.}{Their reflexes are slower.}{Their prefrontal cortex is still maturing while reward-related systems are highly active.}
\end{q}

\section{Gender and sexual orientation (3.3)}

\textbf{Sex} refers to biological characteristics; \textbf{gender} refers to the socially and culturally constructed roles, behaviors, and identities associated with being male or female. \textbf{Gender roles} are expectations for how each gender should act; \textbf{gender typing} is acquiring a traditional role. Gender development is shaped by biology, social learning (modeling and reinforcement), and culture. \textbf{Sexual orientation} is an enduring pattern of attraction; research points to biological contributions (such as prenatal hormone exposure and genetic factors) and finds no evidence that parenting style determines it.

\begin{q}{2}{3.3 Gender}{Learning gender-typed behavior by imitating and being rewarded = social learning.}
A boy is praised for playing with trucks and copies his father's way of fixing things. This best illustrates gender development through
\chh{maturation}{social learning}{critical periods}{assimilation}
\end{q}

\section{Cognitive development (3.4)}

\subsection{Piaget}
Children build \textbf{schemas}. \textbf{Assimilation} fits new information into an existing schema (calling a cow a ``doggy''); \textbf{accommodation} changes the schema to fit (learning ``cow'').
\begin{compare}{Stage (approx.\ age)}{Hallmarks}{Limitations / key test}
Sensorimotor (0--2) & Learning through senses and actions; \textbf{object permanence} develops (around 8 months). & Out of sight, out of mind early on. \\
Preoperational (2--7) & Language, pretend play, symbolic thinking. & \textbf{Egocentrism} (cannot take another's view), no \textbf{conservation}, animism. \\
Concrete operational (7--11) & \textbf{Conservation}, reversibility, logical reasoning about concrete events. & Struggles with abstract hypotheticals. \\
Formal operational (12+) & Abstract, hypothetical, and systematic reasoning. & Not all adults use it consistently. \\
\end{compare}

\subsection{Beyond Piaget}
\textbf{Vygotsky}: development is social and cultural. The \textbf{zone of proximal development} is the gap between what a child can do alone and with help; \textbf{scaffolding} is temporary support that is removed as skill grows. \textbf{Theory of mind}: understanding that others have beliefs and desires different from one's own (typically by about age 4--5; tested with false-belief tasks). Adolescent egocentrism includes the \textbf{imaginary audience} (everyone is watching me) and the \textbf{personal fable} (I am unique and invulnerable).

In adulthood, \textbf{fluid} abilities decline gradually while \textbf{crystallized} knowledge holds or grows; \textbf{neurocognitive disorders} (dementia) are not a normal part of aging.

\begin{q}{2}{3.4 Piaget}{Thinking taller glass has more = lacks conservation (preoperational).}
A 4-year-old insists a tall, thin glass holds more juice than a short, wide glass, even after watching the same juice poured between them. She lacks
\chh{object permanence}{conservation}{a theory of mind}{scaffolding}
\end{q}

\begin{q}{4}{3.4 Vygotsky}{Tasks a child can do with help but not alone fall in the ZPD.}
Ari cannot solve two-step equations alone but can with hints from his tutor, who gradually removes them. The tutor is working within Ari's
\ch{sensorimotor stage}{egocentrism}{concrete operational schema}{zone of proximal development}
\end{q}

\begin{q}{1}{3.4 Piaget}{Fitting new info into an existing schema = assimilation.}
A toddler calls every round object, including an orange, a ``ball.'' She is
\chh{assimilating}{accommodating}{conserving}{scaffolding}
\end{q}

\begin{q}{3}{3.4 Adolescent egocentrism}{Belief one is unique and invulnerable = personal fable.}
A 15-year-old texts while driving, convinced that accidents ``happen to other people.'' This reflects
\chh{the imaginary audience}{theory of mind}{the personal fable}{object permanence}
\end{q}

\section{Communication and language (3.5)}

\begin{terms}
Phoneme & Smallest distinct sound unit (/ch/ in \emph{chat}). \\
Morpheme & Smallest unit that carries meaning (\emph{un-}, \emph{-ed}, \emph{dog}). \\
Semantics & Rules for deriving meaning. \\
Syntax / grammar & Rules for ordering words into sentences; grammar is the whole system of rules. \\
\end{terms}
\textbf{Milestones}: cooing, then \textbf{babbling} (about 4 months; sounds from all languages), \textbf{one-word stage} (about 12 months), \textbf{two-word / telegraphic speech} (about 24 months: ``want milk''), then rapid growth. \textbf{Overgeneralization} (overregularization): applying a grammar rule too broadly (``I goed''). Children exposed to language late in childhood struggle to master grammar, suggesting a \textbf{critical period}. \textbf{Linguistic relativity} (Whorf): language influences thought.

\begin{q}{3}{3.5 Language}{Applying the -ed rule to irregular verbs = overgeneralization.}
A 3-year-old says ``We goed to the park and I runned.'' This shows
\chh{telegraphic speech}{babbling}{overgeneralization}{linguistic relativity}
\end{q}

\begin{q}{2}{3.5 Language}{Smallest meaningful unit = morpheme; ``-s'' marks plural.}
In the word \emph{cats}, the final ``-s'' is best described as a
\chh{phoneme only}{morpheme}{syntax rule}{semantic error}
\end{q}

\section{Social-emotional development (3.6)}

\subsection{Attachment and parenting}
\textbf{Harlow's monkeys} preferred a soft cloth ``mother'' to a wire one that gave milk: attachment rests on \textbf{contact comfort}, not just feeding. \textbf{Lorenz}: goslings \textbf{imprint} on the first moving object during a critical period. \textbf{Temperament} is a baby's inborn emotional reactivity.

\textbf{Ainsworth's Strange Situation} classifies attachment:
\begin{terms}
Secure & Explores with caregiver present; distressed when caregiver leaves; comforted on return. \\
Insecure-avoidant & Little distress at separation; ignores caregiver on return. \\
Insecure-anxious (ambivalent) & Clingy; very distressed; hard to soothe, may resist caregiver on return. \\
Disorganized & Inconsistent, confused, or fearful responses. \\
\end{terms}

\begin{compare}{Parenting style}{Demands / control}{Responsiveness / warmth}
Authoritative & High, with explained rules & High. Linked to the best outcomes (self-reliance, social competence). \\
Authoritarian & High; ``because I said so'' & Low. \\
Permissive & Low & High. \\
Neglectful (uninvolved) & Low & Low. \\
\end{compare}

\subsection{Erikson's psychosocial stages}
\begin{compare}{Life stage}{Crisis}{Positive outcome}
Infancy & Trust vs.\ mistrust & Trust when needs are reliably met. \\
Toddlerhood & Autonomy vs.\ shame and doubt & Independence (``I do it''). \\
Preschool & Initiative vs.\ guilt & Starting tasks and plans. \\
Elementary & Competence (industry) vs.\ inferiority & Pride in accomplishments. \\
Adolescence & Identity vs.\ role confusion & A sense of self. \\
Young adulthood & Intimacy vs.\ isolation & Close relationships. \\
Middle adulthood & Generativity vs.\ stagnation & Contributing to the next generation. \\
Late adulthood & Integrity vs.\ despair & Satisfaction with a life lived. \\
\end{compare}

\subsection{Context}
\textbf{Bronfenbrenner's ecological systems}: microsystem (family, school), mesosystem (links between them), exosystem (parent's workplace), macrosystem (culture), chronosystem (changes over time). \textbf{Adverse childhood experiences} (abuse, neglect, household dysfunction) raise later health risks; \textbf{resilience} protects. \textbf{Emerging adulthood} (about 18 to mid-20s) is an extended period of exploration in many cultures.

\begin{q}{4}{3.6 Attachment}{Distress at departure, comforted on return = secure attachment.}
In the Strange Situation, Mia cries when her mother leaves and quickly calms and resumes playing when she returns. Mia's attachment is
\chh{avoidant}{anxious-ambivalent}{disorganized}{secure}
\end{q}

\begin{q}{1}{3.6 Harlow}{Monkeys chose the cloth mother: contact comfort drives attachment.}
Harlow's research with infant monkeys showed that attachment depends mainly on
\chh{contact comfort}{feeding}{imprinting}{object permanence}
\end{q}

\begin{q}{2}{3.6 Parenting}{High demands with explained reasons and warmth = authoritative.}
Parents who set firm rules, explain the reasons behind them, and are warm and responsive are
\chh{authoritarian}{authoritative}{permissive}{neglectful}
\end{q}

\begin{q}{3}{3.6 Erikson}{Midlife concern for guiding the next generation = generativity vs. stagnation.}
At 48, Rosa begins mentoring young nurses because she wants to give back. In Erikson's theory she is resolving
\ch{intimacy vs.\ isolation}{integrity vs.\ despair}{generativity vs.\ stagnation}{identity vs.\ role confusion}
\end{q}

\section{Classical conditioning (3.7)}

\textbf{Pavlov}: a \textbf{neutral stimulus} (bell) paired with an \textbf{unconditioned stimulus (UCS)} (food) that naturally triggers an \textbf{unconditioned response (UCR)} (salivation) becomes a \textbf{conditioned stimulus (CS)} that triggers a \textbf{conditioned response (CR)} (salivation to the bell).
\begin{terms}
Acquisition & Initial learning of the CS--UCS link; best when the CS comes just before the UCS. \\
Extinction & CR weakens when the CS is presented repeatedly without the UCS. \\
Spontaneous recovery & An extinguished CR reappears after a rest. \\
Generalization & Responding to stimuli similar to the CS. \\
Discrimination & Responding only to the CS, not similar stimuli. \\
Higher-order conditioning & A CS is used to condition a new neutral stimulus. \\
Counterconditioning & Pairing the CS with a new response incompatible with the old one (exposure therapies build on this). \\
Taste aversion & One pairing of a taste with illness creates a lasting aversion, even with hours of delay (Garcia): \textbf{biological preparedness}. \\
Habituation & Decreasing response to a repeated stimulus (a non-associative form of learning). \\
\end{terms}
\textbf{Watson's ``Little Albert''}: pairing a white rat with a loud noise produced fear of the rat that generalized to other furry objects (and raised serious ethical problems).

\begin{trap}
To label a classical conditioning scenario, find the response that happens \emph{naturally} first: whatever caused it is the UCS. The same response to the new cue is the CR. UCR and CR are usually the same behavior.
\end{trap}

\begin{q}{1}{3.7 Classical}{The vaccine shot is the UCS; the clinic smell becomes the CS.}
A child who received painful shots now cries as soon as she smells the clinic's antiseptic. The smell of antiseptic is the
\chh{conditioned stimulus}{unconditioned stimulus}{unconditioned response}{conditioned response}
\end{q}

\begin{q}{4}{3.7 Classical}{Return of an extinguished response after rest = spontaneous recovery.}
A dog's conditioned salivation to a bell was extinguished. After a week away from the lab, the dog salivates again when it hears the bell. This is
\chh{generalization}{acquisition}{discrimination}{spontaneous recovery}
\end{q}

\begin{q}{2}{3.7 Taste aversion}{One-trial learning with a long delay = conditioned taste aversion.}
After getting sick hours after eating sushi, Ben feels nauseated at the sight of sushi for years. This is best explained by
\ch{higher-order conditioning}{conditioned taste aversion}{negative reinforcement}{latent learning}
\end{q}

\section{Operant conditioning (3.8)}

\textbf{Thorndike's law of effect}: behaviors followed by satisfying consequences become more likely. \textbf{Skinner} studied this with operant chambers.

\begin{compare}{}{Add a stimulus (positive)}{Remove a stimulus (negative)}
Behavior \(\uparrow\) (reinforcement) & \textbf{Positive reinforcement}: praise for homework. & \textbf{Negative reinforcement}: buckling the seat belt stops the beeping. \\
Behavior \(\downarrow\) (punishment) & \textbf{Positive punishment}: extra chores for lying. & \textbf{Negative punishment}: losing phone privileges. \\
\end{compare}

\begin{trap}
Negative reinforcement is \emph{not} punishment. ``Negative'' means something is taken away; ``reinforcement'' means the behavior increases. Taking aspirin to remove a headache is negatively reinforced.
\end{trap}

\textbf{Primary reinforcers} satisfy biological needs (food); \textbf{secondary (conditioned) reinforcers} gain value through association (money, grades, tokens). \textbf{Shaping} reinforces successive approximations of a target behavior. \textbf{Instinctive drift}: animals revert to instinctive behaviors that interfere with training. \textbf{Learned helplessness} (Seligman): after repeated uncontrollable bad events, an organism stops trying to escape. \textbf{Superstitious behavior} arises from accidental reinforcement.

\begin{compare}{Schedule}{Reinforcement after}{Example and pattern}
Continuous & Every response & Fastest acquisition; fastest extinction. \\
Fixed ratio (FR) & A set number of responses & Paid per 10 shirts sewn. High rate, pause after reward. \\
Variable ratio (VR) & An unpredictable number of responses & Slot machines. Highest rate, \textbf{most resistant to extinction}. \\
Fixed interval (FI) & First response after a set time & Weekly quiz. ``Scalloped'' pattern: responding speeds up near the deadline. \\
Variable interval (VI) & First response after unpredictable time & Pop quizzes; checking email. Slow, steady rate. \\
\end{compare}

\begin{q}{3}{3.8 Reinforcement}{Removing an aversive stimulus increases the behavior = negative reinforcement.}
Kai cleans his room to make his parents stop nagging him, and he cleans more often in the future. The removal of nagging is
\chh{positive punishment}{negative punishment}{negative reinforcement}{positive reinforcement}
\end{q}

\begin{q}{2}{3.8 Punishment}{Taking away a desired stimulus to decrease behavior = negative punishment.}
A teen loses her car keys for a week after missing curfew, and she stops missing curfew. This is
\chh{positive punishment}{negative punishment}{negative reinforcement}{extinction}
\end{q}

\begin{q}{4}{3.8 Schedules}{Unpredictable number of responses = variable ratio; most resistant to extinction.}
Which schedule describes a fisher who catches a fish after an unpredictable number of casts, and why does she keep casting?
\ch{Fixed interval; she knows when a fish will bite}{Fixed ratio; every tenth cast is rewarded}{Variable interval; time determines the reward}{Variable ratio; the next cast could always pay off, making the behavior resistant to extinction}
\end{q}

\begin{q}{1}{3.8 Shaping}{Reinforcing successive approximations = shaping.}
A trainer rewards a dog first for looking at a hoop, then for walking toward it, then for jumping through it. This is
\chh{shaping}{instinctive drift}{generalization}{learned helplessness}
\end{q}

\section{Social, cognitive, and neurological factors (3.9)}

\begin{terms}
Observational learning & Learning by watching others (\textbf{modeling}). Bandura's \textbf{Bobo doll} study: children who watched an adult hit the doll imitated the aggression. \\
Vicarious conditioning & Learning from the consequences others receive. \\
Insight learning & Sudden realization of a solution (Köhler's chimps stacking boxes). \\
Latent learning & Learning that occurs without reinforcement and appears only when there is an incentive (Tolman's rats formed \textbf{cognitive maps} of a maze). \\
Mirror neurons & Neurons that fire both when performing and when observing an action; proposed basis for imitation and empathy. \\
\end{terms}

\begin{q}{2}{3.9 Latent learning}{Learning without reward revealed once a reward appears = latent learning.}
Rats that explored a maze for days without food suddenly ran it as fast as always-rewarded rats once food was placed at the end. This demonstrates
\ch{insight learning}{latent learning and cognitive maps}{shaping}{spontaneous recovery}
\end{q}

\begin{q}{3}{3.9 Observational}{Learning by watching a model = observational learning (Bandura).}
After watching an older sibling get praised for setting the table, a 5-year-old begins setting the table herself. This best illustrates
\ch{classical conditioning}{insight learning}{observational learning with vicarious reinforcement}{instinctive drift}
\end{q}

\makekey
\end{document}
